\section{Signed Hamiltonian-permanent comparison and parity cancellation}
\label{sec:relativehamilton}
Write \(\operatorname{cyc}D\) for the sum of products over
directed cyclic permutations. For a symmetric kernel of order
at least three it equals \(2Z_{\rm H}(D)\).
The exact identity of Liu~\cite{Liu} is
\[
 N\operatorname{cyc}D
 =\sum_{\emptyset\ne U\subseteq[1,N]}
       |U|(-1)^{|U|-1}\det D[U]\,
                                  \operatorname{per}D[U^c].
\]
We use this signed identity before estimating errors.
Its principal terms must not be replaced by their absolute mass.

The full identity now has the maintained Lean proof
\texttt{hamiltonianLiuIdentity}, over every commutative coefficient ring.
For clarity, the proof expands the actual determinant and permanent and
regroups their products by the resulting full permutation \(\sigma\).
The selected vertex set must be invariant under \(\sigma\), so it is
uniquely a union of complete orbits, including singleton fixed points.
For a selected orbit family \(\mathcal C\), the restricted sign satisfies
\[
 (-1)^{|\bigcup\mathcal C|}
       \operatorname{sgn}(\sigma|_{\bigcup\mathcal C})
                         =(-1)^{|\mathcal C|}.
\]
Consequently the full coefficient of the permutation product is
\[
 -\sum_{\mathcal C\subseteq\operatorname{Orb}(\sigma)}
       (-1)^{|\mathcal C|}|\bigcup\mathcal C|
 =N\,\mathbf1_{\{|\operatorname{Orb}(\sigma)|=1\}}.
\]
Indeed each orbit length is multiplied by the alternating subset sum
over all other orbits. This sum is zero unless there are no other
orbits. The empty selected set contributes zero, including in
dimension zero. The complementary-permutation bijection and its
matrix-product identity are part of the formal proof, so no fixed-point,
overlap or completion term is omitted. This formalizes the established
identity; it is not a new enumeration claim.

\paragraph{Formalized weighted row-cap permanent bound.}
This bound is used only for a separately retained tail.
Let \(M\) be nonnegative of order \(n\geq2\), with row sums at
most \(1\) and entries at most \(c>0\).
Write \(w_\sigma=\prod_iM_{i,\sigma(i)}\) and
\(P=\sum_\sigma w_\sigma\).
If \(P=0\), every upper bound below is immediate.
For each fixed row order, weighted sequential revelation gives
a subprobability \(q_\sigma\), with its actual available-column
denominators. If \(C_\sigma\) is the sum of their logarithms,
the inequality \(\log x\leq x-1\) gives
\[
 \log P\leq P^{-1}\sum_\sigma w_\sigma C_\sigma.
\]
Zero matching weights contribute zero; no positive-entry
approximation is needed.

Fix a matching of positive weight and average over all row orders.
When row \(i\) is revealed with \(k\) other rows remaining,
the latter are a uniform subset of the other \(n-1\) rows.
The available mass therefore has conditional mean
\[
 M_{i,\sigma(i)}
       +\frac{k}{n-1}\sum_{j\ne i}M_{i,\sigma(j)}
                         \leq c+\frac{k}{n-1}.
\]
Applying the logarithmic mean inequality at each stage, then
averaging with the full matching weights, proves
\[
 \operatorname{per}M\leq
 \exp\!\left(\sum_{k=0}^{n-1}\log\left(c+\frac{k}{n-1}\right)\right).
\]
For \(m=n-1\), the logarithm is at most
\[
 \log c-m+\log m+1+cm(1+\log m),
\]
by \(\log(1+u)\leq u\), the harmonic bound and the Stirling
upper bound for the factorial logarithm.
Consequently, if \(A\geq1\), \(N\geq1\), \(n\leq N\) and
\(c=A/N\), then
\[
 \operatorname{per}M\leq \exp(-n+16A\log(2N)).
\]
The same conclusion for \(n=0,1\) follows directly.
The weighted argument, explicit constant, exceptional orders and
transfer to every actual independent-support rectangular minor
have maintained Lean proofs. In particular, for a doubly stochastic
matrix \(D\) with maximum entry \(A/N\),
\[
 \operatorname{per}D[I^c,J^c]
                       \leq \exp(r-N+16A\log(2N)),\qquad |J|=r.
\]
The denominator bound \(N!/N^N\geq e^{-N}\), together with
van der Waerden, gives the ordinary-proof ratio estimate
\[
 \frac{\operatorname{per}D[I^c,J^c]}{\operatorname{per}D}
                       \leq\exp(r+16A\log(2N)).
\]
For the admissible stochastic expander kernels at sufficiently large
order, the maintained proof instead consumes the already formalized
normalized positivity data.  It proves
\(\operatorname{per}D\geq e^{-N}/2\), hence the all-rank bound
\[
 0\leq\frac{\operatorname{per}D[I^c,J^c]}{\operatorname{per}D}
                       \leq2\exp(r+16A\log(2N)).
\]
The factor \(2\) is retained throughout the full compensated bipartite
tail.  Its tilted Cauchy--Binet summation, including the constant and
first-order subtractions, now has a maintained proof and gives
\(\exp(-\log(2N)^3\log\log(2N)/4)\) above one uniform threshold
when \(A\leq C_0\log(2N)\).  The proof and its explicit scalar budget
are given in the signed comparison section.

\begin{theorem}[A relative comparison with a growing density parameter]
\label{thm:relativeaperiodic}
Fix \(\delta\in(0,1)\). Let \(D_N\) be real nonnegative
doubly stochastic matrices, and put
\[
 B=D_N-J_N/N,\qquad \|B\|\leq1-\delta,\qquad
 \max D_{ij}\leq A_N/N,\qquad A_N\geq2.
\]
If \(A_N^2=o(\log N)\), then
\[
 \operatorname{cyc}D_N
 =(1+o(1))\frac{(N-1)!}{N^N}
       \frac{\det(I-B)}{\sqrt{\det(I-B^\top B)}}.
\]
The relative error is
\(O_\delta(N^{-1/9}\exp(C_\delta A_N^2))\), and hence
tends to zero along every sequence with the stated entry
growth and fixed gap. No positive-entry lower bound is assumed.
\end{theorem}
\begin{proof}
For all sufficiently large \(N\), \(A_N\leq\log N\).
The permanent and compensated-minor estimates of
Section~\ref{sec:relativepermanents} therefore apply
with a fixed polylogarithmic exponent.
Put \(L=\log(2N)\), and split Liu's sum at \(r=L^3\),
with integer rounding at this cutoff.

For \(r\leq L^3\), the minor ratio is
\(1+O_\delta(N^{-1/9})\), uniformly in \(U\).
Hadamard's inequality gives
\[
 \sum_{|U|=r}|\det D[U]|
   \leq\binom Nr\left(\frac{A\sqrt r}{N}\right)^r
   \leq(eA/\sqrt r)^r .
\]
Consequently
\[
 \sum_{r\geq1}r\sum_{|U|=r}|\det D[U]|
                         \leq\exp(C A^2).
\]
For example split at a sufficiently large constant times
\(A^2\); before that cutoff each logarithmic summand is
\(O(A^2)\), and afterwards the series decreases geometrically.
The small-minor error is therefore at most
\(N^{-1/9}\exp(C A^2)\).

For \(r>L^3\), the separately retained permanent-ratio bound
above yields total absolute tail at most
\[
 \exp(C A L)\sum_{r>L^3}
                    r(e^2A/\sqrt r)^r
       \leq\exp(-c L^3\log L),
\]
for large \(N\), since \(A\leq\log N\).
The same estimate without the permanent ratio controls the
tail of the determinant-only sum.
Thus the entire signed leading expression can be evaluated:
\[
 \frac{N\operatorname{cyc}D}{\operatorname{per}D}
 =\sum_{\emptyset\ne U}|U|(-1)^{|U|-1}\det D[U]
    +O_\delta\left(N^{-1/9}e^{CA^2}
                         +e^{-cL^3\log L}\right)
 =\det(I-B)+\text{the stated error}.
\]
For the last equality,
\(\det(I-tD)=(1-t)\det(I-tB)\), using both stochastic row
and column equations; differentiate at \(t=1\).
No determinant contribution has been dropped termwise.

We check the relative normalization.
The convergent trace expansion gives
\[
 \log\det(I-B)
       =-\operatorname{tr}B-\sum_{j\geq2}
                               \operatorname{tr}(B^j)/j .
\]
Here \(|\operatorname{tr}B|\leq A+1\), and
\[
 |\operatorname{tr}(B^j)|
       \leq\|B\|^{j-2}\|B\|_F^2
       \leq(1-\delta)^{j-2}A .
\]
Therefore \(\det(I-B)\geq\exp(-C_\delta A)>0\).
The normalized error is at most
\[
 O_\delta\left(
 N^{-1/9}\exp(C A^2+C_\delta A)
              +\exp(-cL^3\log L+C_\delta A)\right)=o(1).
\]
Apply Theorem~\ref{thm:growingpermanent} to the remaining full
permanent. Its determinant and factorial factors give the
displayed formula.
This extends the fixed-entry conclusion
of~\cite[Proposition 5.1]{DW} by an explicit growing-parameter
argument; it does not claim the underlying identity as new.
\end{proof}

\begin{theorem}[Relative bipartite universality]
\label{thm:relativebipartite}
Let \(T=T_m\) be real nonnegative doubly stochastic, with
\[
 B=T-J_m/m,\qquad \|B\|\leq1-\delta,\qquad
 \max T_{ij}\leq A_m/m,\qquad 2\leq A_m=o(\log m),
\]
for fixed \(\delta\in(0,1)\).
For the symmetric bipartite kernel
\(P=\begin{pmatrix}0&T\\T^\top&0\end{pmatrix}\),
\[
 Z_{\rm H}(P)
 =(1+o(1))\frac{(m!)^2}{2m\,m^{2m}}
 =(1+o(1))\pi e^{-2m}.
\]
The relative error is
\(O_\delta(m^{-1/9}\exp(C_\delta A_m))\).
In particular the leading relative constant is independent
of the nontrivial singular spectrum of \(T\).
The exact parity eigenvalue is treated separately, not inserted
into a formula requiring norm strictly below one.
\end{theorem}
\begin{proof}
For \(m\geq2\), Liu's identity for \(P\) has contributions
only when the selected set contains equally many vertices
from the two parts.
Writing those subsets as \(I,J\), each of size \(r\), exactly
\[
 \det P[I\cup J]=(-1)^r\det T[I,J]^2,\qquad
 \operatorname{per}P[(I\cup J)^c]
                         =\operatorname{per}T[I^c,J^c]^2.
\]
Consequently the complete signed expression is
\[
 2m Z_{\rm H}(P)
 =\sum_{r=1}^m(-1)^{r+1}r
     \sum_{\substack{|I|=|J|=r}}
       \det T[I,J]^2\,\operatorname{per}T[I^c,J^c]^2.
\]
This retains every mismatched row/column support and every
complementary permanent.
Its directed form is now the maintained theorem
\texttt{hamiltonianBipartiteCycleIdentity}:
\(m\operatorname{cyc}P\) equals the complete signed sum on the right.
The displayed unoriented normalization follows from
\(\operatorname{cyc}P=2Z_{\rm H}(P)\) for symmetric order \(2m\geq4\).
The formal proof transports both selected supports and both complements
explicitly; their equal-cardinality enumerations are independent.
Unbalanced determinant terms are proved zero, and the balanced
determinant and permanent squares are actual matrix identities.
The empty supports remain in the code with their zero cardinality
coefficient. No minor-ratio estimate is part of this exact identity.

The compensated-minor lemma gives a uniform
\(1+O_\delta(m^{-1/9})\) ratio for \(r\leq(\log(2m))^3\).
By Cauchy--Binet,
\[
 \sum_{|I|=|J|=r}\det T[I,J]^2=e_r(TT^\top),
\]
where \(e_r\) is the elementary symmetric polynomial in the
eigenvalues.
The total absolute mass for the small-minor error is now
bounded more sharply than the general determinant bound:
\[
 \sum_r r e_r(TT^\top)
 \leq\operatorname{tr}(TT^\top)\exp(\operatorname{tr}(TT^\top))
 \leq A e^A.
\]
The last inequality uses stochasticity,
\(\sum_{ij}T_{ij}^2\leq m\max T_{ij}\leq A\).

The retained large-minor tail is also explicit.
Hadamard and the two subset counts give
\[
 \sum_{|I|=|J|=r}\det T[I,J]^2
          \leq(e^2A^2/r)^r .
\]
Together with the squared permanent-ratio bound, its normalized
tail for \(r>(\log(2m))^3\) is at most
\[
 e^{C A\log(2m)}
       \sum_{r>(\log(2m))^3}r(e^4A^2/r)^r
          \leq e^{-c(\log(2m))^3\log\log(2m)} .
\]
Here \(A\leq\log m\) eventually.
The determinant-only tail satisfies the same smaller bound.

Now evaluate the signed leading sum, rather than bounding
it by \(Ae^A\).
Since \(TT^\top=J_m/m+BB^\top\),
\[
 \det(I-xTT^\top)=(1-x)\det(I-xBB^\top).
\]
Its derivative at \(x=1\) proves the exact cancellation
\[
 \sum_r(-1)^{r+1}r e_r(TT^\top)=\det(I-BB^\top).
\]
Thus
\[
 \frac{2m Z_{\rm H}(P)}{(\operatorname{per}T)^2}
 =\det(I-BB^\top)
      +O_\delta(m^{-1/9}Ae^A)
      +e^{-c(\log(2m))^3\log\log(2m)}.
\]
The determinant is at least \(e^{-C_\delta A}\).
The relative error therefore tends to zero because
\(A=o(\log m)\).
Finally
\[
 (\operatorname{per}T)^2
  =(1+o(1))\left(\frac{m!}{m^m}\right)^2
                               \det(I-B^\top B)^{-1}.
\]
The two determinants are equal.
Their entire spectral contribution cancels, giving the
factorial formula and then its Stirling form.
The argument shows why an absolute estimate of the signed
leading minors would miss the relative constant.
\end{proof}

\begin{corollary}[Relative counts from entropy-normalized graph weights]
\label{cor:relativegraphentropy}
For a simple balanced bipartite graph with positive stochastic
scaling \(P=\diag(u)A_G\diag(u)\), assume the hypotheses of
Theorem~\ref{thm:relativebipartite}.
Then
\[
 h(G)=(1+o(1))\pi
       \exp\left(-N-2\sum_v\log u_v\right).
\]
For a simple graph of either parity with stochastic scaling
satisfying Theorem~\ref{thm:relativeaperiodic}, its relative
count is
\[
 h(G)=(1+o(1))\frac{(N-1)!}{2N^N}
  \exp\left(-2\sum_v\log u_v\right)
           \sqrt{\frac{\det(I-B)}{\det(I+B)}} ,
 \qquad B=P-J_N/N .
\]
Both formulas allow zero support. The second uses symmetry
to simplify the determinant ratio.
\end{corollary}
\begin{proof}
Every cycle weight is multiplied by \(\prod_vu_v^2\).
For symmetric \(B\),
\(I-B^\top B=(I-B)(I+B)\), and both factors are positive
definite on their nontrivial spaces under the stated gap.
Substitute the two relative theorems, dividing directed
cycle counts by two only in the aperiodic formula.
\end{proof}

\begin{remark}[The density ranges are part of the result]
The relative aperiodic estimate here requires \(A^2=o(\log N)\);
the bipartite estimate permits \(A=o(\log N)\).
They are stronger precision statements on a narrower domain
than the earlier diffuse-kernel entropy theorem.
These estimates alone do not apply to the prime-sum
normalization where \(N\max P_{ij}\) is of order \(\log N\).
The signed first-correction argument in
Section~\ref{sec:signedjetcomparison} closes that boundary
under its explicit determinant-amplification condition.
They do cover all unitary Cayley graphs, whose corresponding
parameter is \(O(\log\log N)\), as proved next.
\end{remark}
